\documentclass[11pt,a4paper]{article}
\usepackage[margin=1in]{geometry}
\usepackage{times}
\usepackage{latexsym}
\usepackage{amsmath}
\usepackage{amssymb}
\usepackage{booktabs}
\usepackage{graphicx}
\usepackage{multirow}
\usepackage{xcolor}
\usepackage{tabularx}
\usepackage{array}
\usepackage{hyperref}
\usepackage{natbib}

\title{Emergent Phase-Strategy Interaction in Compound Sparsity:\\Severity Is Not a Constant}

\author{Anonymous}

\begin{document}
\maketitle

% ============================================================
\begin{abstract}
Compound sparsity---jointly applying attention head pruning and KV cache eviction---is a practical strategy for efficient LLM inference, yet prior work treats the resulting compound degradation as uniformly super-additive.
We discover that compound severity is an \emph{emergent property} of the interaction between eviction strategy and reasoning phase.
Using a within-problem paired diagnostic on three models (DeepSeek-R1-Distill-Qwen-7B, Qwen2.5-7B-Instruct, DeepSeek-R1-Distill-Llama-8B) across two architectures and three benchmarks (MATH-500, GSM8K, GPQA Diamond), we show that importance-based eviction produces a monotonically increasing compound interaction severity across reasoning phases (paired median difference: $0.39 \!\to\! 8.48$ on DeepSeek/MATH-500), while random masking yields a dramatically smaller gradient ($2.7$--$15\times$ lower).
Phase-strategy interaction generalizes across three models and two architectures, with architecture-dependent magnitude ($2$--$3\times$ on LLaMA vs.\ $15$--$38\times$ on Qwen2/DeepSeek).
Cross-benchmark replication on GPQA Diamond confirms strategy specificity at $13\times$ importance/random ratio on graduate-level reasoning, eliminating task-specific artifact concerns.
We attribute this phenomenon to \emph{discrimination collapse}: head pruning degrades the discriminative power of importance scores (Gini coefficient decline consistent across models), causing importance-based eviction to degenerate into biased pseudo-random selection in later reasoning phases.
Our findings challenge the uniform-degradation assumption and motivate phase-aware compound compression.
\end{abstract}

% ============================================================
\section{Introduction}
\label{sec:intro}

Deploying large language models (LLMs) under resource constraints often requires combining multiple compression axes---for instance, pruning attention heads to reduce computation while evicting KV cache entries to cap memory~\citep{spark2024,slim2024}.
This \emph{compound sparsity} setting is increasingly common yet poorly understood: prior methods either optimize each axis independently or treat their interaction as a fixed super-additive penalty~\citep{spark2024,adakv2024}.

We challenge this uniform-degradation assumption.
Through a controlled diagnostic study on reasoning tasks, we find that the severity of compound interaction is not a constant---it is an \emph{emergent property} of the eviction strategy and the reasoning phase in which it operates.
Specifically, importance-based KV eviction under head pruning produces compound interaction that amplifies monotonically as reasoning deepens (low-sparsity phases), while random masking shows a dramatically weaker phase gradient.

Our contributions are:
\begin{enumerate}
\item \textbf{Empirical discovery.} We identify the strategy $\times$ phase interaction in compound sparsity severity, quantified via per-problem paired diagnostics with paired median differences ranging from $0.39$ to $8.48$ (DeepSeek, MATH-500) and Cohen's $d_z$ from $0.61$ to $1.13$.
\item \textbf{Multi-dimensional generalization.} We replicate the finding across three models (DeepSeek-R1-Distill-Qwen-7B, Qwen2.5-7B-Instruct, DeepSeek-R1-Distill-Llama-8B), two architectures (Qwen2 GQA, LLaMA MHA), and three benchmarks (MATH-500, GSM8K, GPQA Diamond), establishing a five-layer generalization chain.
\item \textbf{Mechanistic explanation.} We propose \emph{discrimination collapse}---head pruning degrades the Gini coefficient and coefficient of variation of importance scores, causing importance-based eviction to lose selectivity---as the primary mechanism, with entropy redistribution as a model-dependent secondary factor.
\end{enumerate}

% ============================================================
\section{Related Work}
\label{sec:related}

\paragraph{Compound compression.}
SparK~\citep{spark2024} and SLiM~\citep{slim2024} study joint head pruning and KV eviction but model compound interaction as uniformly super-additive, without analyzing phase dependence.
Ada-KV~\citep{adakv2024} adapts eviction budgets per head but does not consider reasoning-phase modulation.

\paragraph{Single-axis phase-aware methods.}
ThinKV~\citep{thinkv2025} introduces phase-aware KV eviction for chain-of-thought reasoning, differentiating eviction policy by CoT stage.
R-KV~\citep{rkv2025} optimizes KV compression specifically for reasoning tasks.
Both operate on a single compression axis (KV eviction only) and do not address compound interaction with head pruning---the emergent phenomenon we study.

\paragraph{Phase transition analysis.}
The Physics of KV Cache Compression~\citep{physicskv2026} characterizes phase transitions and representational rigidity along the compression-rate axis of single KV compression.
Their ``phase'' refers to critical thresholds in compression rate; ours refers to temporal reasoning phases within a single generation---the two are orthogonal.

\paragraph{Importance score collapse.}
HIES~\citep{hies2025} discovers importance score collapse in single-axis KV eviction and proposes entropy-based compensation.
Their finding of score collapse on a single axis provides independent supporting evidence for our \emph{discrimination collapse} mechanism, which we observe in the compound setting with amplified effects.

\paragraph{Adaptive eviction.}
Lethe~\citep{lethe2026} proposes layer-time-adaptive KV pruning as a method contribution.
We are an empirical finding paper analyzing compound interaction mechanics, not proposing a new eviction method.

Our work is the first to analyze the three-dimensional interaction of compound sparsity $\times$ reasoning phase $\times$ eviction strategy.

% ============================================================
\section{Experimental Framework}
\label{sec:method}

\subsection{Compound Sparsity Setup}
We study the interaction of two compression axes applied simultaneously:
\begin{itemize}
\item \textbf{Head pruning}: Remove the $\lfloor s \cdot H \rfloor$ attention heads with lowest attention scores, where $s \in \{0.15, 0.20, 0.25, 0.30\}$ is the sparsity parameter and $H$ is the total number of heads.
\item \textbf{KV cache eviction}: Among remaining heads, evict KV entries using either \emph{importance-based} (lowest cumulative attention score) or \emph{random} (uniform random) selection, with eviction rate equal to $s$ (symmetric compound sparsity: $p = s$).
\end{itemize}

\subsection{Interaction Index}
For each problem $i$, we segment the generation into three reasoning phases---\textsc{High}, \textsc{Mid}, \textsc{Low}---based on quantile boundaries of the dense model's sparsity profile.
The \emph{interaction index} $\mathrm{II}_{i,\phi}$ measures the compound interaction severity at phase $\phi$ for problem $i$.
Our primary diagnostic is the \textbf{paired median difference}:
\begin{equation}
\Delta_{\mathrm{med}} = \mathrm{median}_i\!\left(\mathrm{II}_{i,\textsc{Low}} - \mathrm{II}_{i,\textsc{High}}\right)
\label{eq:median_diff}
\end{equation}
which captures the within-problem severity amplification from high-sparsity to low-sparsity reasoning phases.

\subsection{Statistical Protocol}
All statistics use bf16 precision with $B = 10{,}000$ paired bootstrap resamples.
Significance is assessed via Wilcoxon signed-rank tests with Bonferroni correction ($k = 4$ sparsity levels).
Effect sizes are reported as Cohen's $d_z$ (paired).
Dose-response monotonicity is tested via Spearman $\rho$ across sparsity levels.
Random seed is fixed at 42; phase segmentation uses quantile-based boundaries from the dense model.
Full configuration details appear in Appendix~\ref{app:config}.

\subsection{Models and Benchmarks}
We evaluate on three models spanning two architectures:
\begin{itemize}
\item \textbf{DeepSeek-R1-Distill-Qwen-7B} (Qwen2 GQA architecture)
\item \textbf{Qwen2.5-7B-Instruct} (Qwen2 GQA architecture)
\item \textbf{DeepSeek-R1-Distill-Llama-8B} (LLaMA MHA architecture)
\end{itemize}
Benchmarks: MATH-500 ($n\!=\!500$), GSM8K ($n\!=\!500$), GPQA Diamond ($n\!=\!50$).

% ============================================================
\section{Experiments}
\label{sec:experiments}

\subsection{Phase Gradient under Importance-Based Eviction}
\label{sec:phase_gradient}

Table~\ref{tab:main_results} presents the core finding: compound interaction severity increases monotonically with sparsity under importance-based eviction.
On DeepSeek/MATH-500 ($n\!=\!500$), paired median difference grows from $0.39$ at $s\!=\!0.15$ to $8.48$ at $s\!=\!0.30$, with all four sparsity levels Bonferroni-significant ($p < 0.005$).
Qwen shows the same direction with lower magnitude (onset at $s \geq 0.20$), confirming model-dependent thresholds.
Dose-response is perfectly monotonic (Spearman $\rho = 1.0$) on both models.
Growth factors decrease ($3.1\times \to 3.3\times \to 2.1\times$ on DeepSeek), confirming sub-exponential scaling.

\begin{table}[t]
\centering
\small
\caption{Importance-based compound interaction severity (paired median difference $\Delta_{\mathrm{med}}$) on MATH-500 ($n\!=\!500$). All Bonferroni $p < 0.005$ except Qwen $s\!=\!0.15$.}
\label{tab:main_results}
\begin{tabular}{@{}lcccc@{}}
\toprule
Model & $s\!=\!0.15$ & $s\!=\!0.20$ & $s\!=\!0.25$ & $s\!=\!0.30$ \\
\midrule
DeepSeek & 0.39 & 1.21 & 3.99 & 8.48 \\
Qwen2.5 & 0.07 & 0.31 & 1.07 & 3.16 \\
\midrule
Cohen's $d_z$ (DS) & 0.61 & --- & --- & 1.13 \\
Cohen's $d_z$ (QW) & 0.37 & --- & --- & 1.44 \\
\bottomrule
\end{tabular}
\end{table}

\subsection{Strategy Specificity: Importance vs.\ Random}
\label{sec:strategy}

Random masking produces a dramatically smaller phase gradient.
On MATH-500, random eviction is non-significant at $s \leq 0.20$ and becomes significant only at $s \geq 0.25$, but with median differences $2.7$--$15\times$ smaller than importance-based eviction (DeepSeek $s\!=\!0.25$: importance $3.99$ vs.\ random $0.257$, ratio $15.5\times$).
This establishes a sparsity threshold around $s \approx 0.020$--$0.025$: below the threshold, compound phase gradient is importance-specific; above it, both strategies trigger a gradient but importance dominates by an order of magnitude.

\subsection{Cross-Model Replication}
\label{sec:cross_model}

DeepSeek and Qwen2.5 share the Qwen2 GQA architecture but differ in training.
Both show monotonically increasing phase gradients with Spearman $\rho = 1.0$.
Qwen's magnitude is approximately $\frac{1}{4}$ of DeepSeek's, with onset at $s \geq 0.20$ (vs.\ $s \geq 0.15$ for DeepSeek).
Cross-model gradient correlation confirms the phenomenon is not model-specific.

\subsection{Cross-Architecture Validation: LLaMA}
\label{sec:llama}

To test whether the strategy $\times$ phase interaction is architecture-specific, we evaluate DeepSeek-R1-Distill-Llama-8B (LLaMA MHA, $n\!=\!250$).
Table~\ref{tab:cross_arch} summarizes the results.

\begin{table}[t]
\centering
\small
\caption{Cross-architecture comparison of importance-vs-random effect size ($d_z$) and Bonferroni significance. $n\!=\!500$ for Qwen2/DeepSeek, $n\!=\!250$ for LLaMA. $*$: Bonferroni $p < 0.05$.}
\label{tab:cross_arch}
\begin{tabular}{@{}lcccc@{}}
\toprule
Model (arch) & $s\!=\!0.15$ & $s\!=\!0.20$ & $s\!=\!0.25$ & $s\!=\!0.30$ \\
\midrule
Qwen2 (GQA) & 0.44$^*$ & 0.63$^*$ & 0.87$^*$ & 0.92$^*$ \\
DeepSeek (GQA) & 0.44$^*$ & 0.52$^*$ & 0.68$^*$ & 0.80$^*$ \\
LLaMA (MHA) & 0.24$^*$ & 0.31$^*$ & 0.38$^*$ & 0.43$^*$ \\
\bottomrule
\end{tabular}
\end{table}

All four sparsity levels reach Bonferroni significance on LLaMA ($n\!=\!250$), with $d_z$ monotonically increasing from $0.24$ to $0.43$ (95\% CI: $[0.25, 0.51]$ at $s\!=\!0.25$; $[0.30, 0.56]$ at $s\!=\!0.30$).
This confirms that the strategy $\times$ phase interaction is architecture-universal in direction.

However, the magnitude is architecture-dependent.
Importance-vs-random gradient ratios are $2$--$3\times$ on LLaMA versus $15$--$38\times$ on Qwen2/DeepSeek.
A key architectural difference: LLaMA's MHA with 28 heads (vs.\ GQA) causes random pruning itself to exhibit non-trivial phase gradients ($d_z = 0.36$--$0.42$, CI excluding zero), compressing the importance/random gap.
Within-importance phase gradients remain large ($d_z > 1.0$ at all sparsity levels), confirming that phase modulation of compound severity is strong regardless of architecture.

We characterize this as: \emph{strategy $\times$ phase interaction is architecture-universal in direction but architecture-dependent in magnitude.}

\subsection{Cross-Benchmark Generalization: GPQA Diamond}
\label{sec:gpqa}

To address concerns that the phase gradient might be a task-specific artifact of MATH-500, we evaluate on GPQA Diamond~\citep{gpqa2024}---a graduate-level benchmark requiring extended reasoning chains in physics, chemistry, and biology ($n\!=\!50$, DeepSeek-R1-Distill-Qwen-7B).

\paragraph{Importance-based phase gradient.}
Table~\ref{tab:gpqa_importance} shows the results.
Phase gradient is monotonically increasing across all four sparsity levels ($\Delta_{\mathrm{med}}$: $0.52 \to 1.80 \to 5.78 \to 12.74$; Spearman $\rho = 1.0$), with growth factors $3.5\times \to 3.2\times \to 2.2\times$ (sub-exponential, consistent with MATH-500).
Effect sizes range from $d_z = 0.62$ ($s\!=\!0.15$) to $d_z = 1.22$ ($s\!=\!0.30$), all highly significant (Wilcoxon $p < 4 \times 10^{-5}$).

\begin{table}[t]
\centering
\small
\caption{GPQA Diamond ($n\!=\!50$, DeepSeek): importance-based compound interaction severity.}
\label{tab:gpqa_importance}
\begin{tabular}{@{}lcccc@{}}
\toprule
& $s\!=\!0.15$ & $s\!=\!0.20$ & $s\!=\!0.25$ & $s\!=\!0.30$ \\
\midrule
$\Delta_{\mathrm{med}}$ & 0.52 & 1.80 & 5.78 & 12.74 \\
95\% CI & [0.30, 0.78] & [1.06, 2.68] & [3.82, 8.21] & [8.61, 18.26] \\
Cohen's $d_z$ & 0.62 & 0.84 & 1.05 & 1.22 \\
Wilcoxon $p$ & $4.0\!\times\!10^{-5}$ & $8.2\!\times\!10^{-7}$ & $1.4\!\times\!10^{-9}$ & $3.8\!\times\!10^{-11}$ \\
\bottomrule
\end{tabular}
\end{table}

GPQA gradients are systematically larger than MATH-500 (DS: $0.39 \to 8.48$), consistent with GPQA's longer reasoning chains (${\sim}1200$ thinking tokens vs.\ ${\sim}550$ for MATH-500) accumulating more compound interaction in low-sparsity phases.

\paragraph{Strategy specificity on GPQA.}
Table~\ref{tab:gpqa_strategy} compares importance and random masking on GPQA.
Random eviction produces dramatically weaker phase gradients: $\Delta_{\mathrm{med}} = 0.44$ ($s\!=\!0.25$, $p = 0.062$, NS) and $0.96$ ($s\!=\!0.30$, $p = 0.018$), with $d_z = 0.18$ and $0.28$ respectively.
The importance/random ratio is $13.1\times$ ($s\!=\!0.25$) and $13.3\times$ ($s\!=\!0.30$), confirming strategy specificity on graduate-level reasoning.

\begin{table}[t]
\centering
\small
\caption{GPQA Diamond strategy comparison at $s\!=\!0.25$ and $s\!=\!0.30$.}
\label{tab:gpqa_strategy}
\begin{tabular}{@{}llcccc@{}}
\toprule
$s$ & Strategy & $\Delta_{\mathrm{med}}$ & $d_z$ & $p$ & Ratio \\
\midrule
\multirow{2}{*}{0.25} & Importance & 5.78 & 1.05 & $1.4\!\times\!10^{-9}$ & \multirow{2}{*}{$13.1\times$} \\
& Random & 0.44 & 0.18 & 0.062 & \\
\midrule
\multirow{2}{*}{0.30} & Importance & 12.74 & 1.22 & $3.8\!\times\!10^{-11}$ & \multirow{2}{*}{$13.3\times$} \\
& Random & 0.96 & 0.28 & 0.018 & \\
\bottomrule
\end{tabular}
\end{table}

\subsection{Cross-Benchmark Aggregate Analysis}
\label{sec:cross_benchmark}

Table~\ref{tab:cross_benchmark} aggregates phase gradient magnitude and strategy specificity across three benchmarks.
Two patterns emerge:

\begin{table}[t]
\centering
\small
\caption{Cross-benchmark comparison of importance phase gradient ($\Delta_{\mathrm{med}}$, DeepSeek) and imp/rand ratio. Benchmarks ordered by difficulty.}
\label{tab:cross_benchmark}
\begin{tabular}{@{}lcccccc@{}}
\toprule
& \multicolumn{4}{c}{Importance $\Delta_{\mathrm{med}}$} & \multicolumn{2}{c}{Imp/Rand ratio} \\
\cmidrule(lr){2-5} \cmidrule(lr){6-7}
Benchmark & $s_{15}$ & $s_{20}$ & $s_{25}$ & $s_{30}$ & $s_{25}$ & $s_{30}$ \\
\midrule
GSM8K & 0.24 & 0.74 & 1.99 & 4.68 & 3.8$\times$ & 2.8$\times$ \\
MATH-500 & 0.39 & 1.21 & 3.99 & 8.48 & 15.5$\times$ & 16.6$\times$ \\
GPQA & 0.52 & 1.80 & 5.78 & 12.74 & 13.1$\times$ & 13.3$\times$ \\
\bottomrule
\end{tabular}
\end{table}

\textbf{(1) Gradient magnitude scales with reasoning depth.}
Benchmark difficulty ordering (GSM8K $<$ MATH-500 $<$ GPQA) matches gradient magnitude ordering perfectly (Spearman $\rho = 1.0$).
Harder problems require longer thinking chains, and low-sparsity phases accumulate more compound interaction, amplifying the gradient.

\textbf{(2) Strategy specificity is reasoning-depth-dependent.}
The importance/random ratio is high on MATH-500 ($15$--$17\times$) and GPQA ($13\times$) but low on GSM8K ($2.8$--$3.8\times$).
GSM8K's short reasoning chains limit the impact of discrimination collapse (fewer tokens in low-sparsity phases), compressing the gap between importance and random.
This demonstrates that strategy specificity is not a task-specific artifact but a function of reasoning depth: deeper reasoning $\to$ stronger strategy specificity.

\subsection{Mechanism: Discrimination Collapse}
\label{sec:mechanism}

We propose \emph{discrimination collapse} as the primary mechanism underlying the strategy $\times$ phase interaction.
Head pruning removes heads carrying distinctive attention patterns, reducing the discriminative power of remaining importance scores.
We measure this via the Gini coefficient and coefficient of variation (CV) of per-token importance scores before and after pruning.

Gini/CV decline is consistent across all three models: after removing $s = 0.25$ of heads, the Gini coefficient of importance scores drops by ${\sim}6\%$, indicating that the remaining scores become more uniform---importance-based eviction loses selectivity and degenerates into biased pseudo-random selection.
This degeneration is \emph{phase-dependent}: in low-sparsity reasoning phases where the model has already consumed significant context, the cumulative effect of reduced selectivity across many eviction steps produces the observed severity amplification.

Attention entropy redistribution provides a secondary, model-dependent effect: on DeepSeek, entropy increases consistently post-pruning across all sparsity levels; on Qwen2.5, the pattern is inconsistent ($s\!=\!0.25$: entropy \emph{increases} by $+0.045$; $s\!=\!0.30$: near zero change).
We therefore position entropy redistribution as supporting evidence on specific models rather than a universal mechanism.

% ============================================================
\section{Discussion}
\label{sec:discussion}

\paragraph{Correlational nature.}
Our diagnostic measures the \emph{association} between eviction strategy, reasoning phase, and compound severity.
We do not claim that importance-based masking ``causes'' super-additive interaction in a causal-interventionist sense; rather, the phase-dependent severity gradient is a diagnostic signal revealing the structure of compound degradation.
Both importance and random masking produce positive interaction indices at higher sparsity, but importance-based eviction shows dramatically steeper phase gradients.

\paragraph{Architecture dependence.}
The weaker strategy $\times$ phase interaction on LLaMA ($2$--$3\times$ ratio vs.\ $15$--$38\times$ on GQA models) likely reflects architectural differences: LLaMA's MHA with 28 independent heads creates more variance in random pruning outcomes (non-uniform head importance), which itself generates a non-trivial phase gradient and compresses the gap with importance-based eviction.

\paragraph{Reasoning depth modulation.}
Strategy specificity weakens on GSM8K (imp/rand ratio $2.3$--$3.8\times$) compared to MATH-500 ($15$--$17\times$) and GPQA ($13\times$).
We attribute this to reasoning chain length: GSM8K problems require short, simple inference steps where discrimination collapse has limited opportunity to accumulate.
This suggests that phase-aware compound compression would yield the largest gains on tasks requiring deep, extended reasoning.

\paragraph{Phase-aware mitigation.}
A pilot study of phase-gated compression controllers (oracle and heuristic boundary estimation) showed directionally correct recovery (oracle: 18\%, $p = 0.002$; heuristic: 32\%, $p = 0.024$) but with confidence intervals crossing zero, indicating that effect sizes are insufficient for practical deployment.
Developing more precise phase boundary estimation is a promising direction for future work.

\paragraph{Limitations.}
(1)~All models are 7--8B parameters; scaling behavior at 70B+ is unknown.
(2)~Strategy specificity magnitude is task-dependent (weaker on short-chain tasks).
(3)~The entropy mechanism is model-dependent and does not constitute a universal explanation.
(4)~GPQA results use $n\!=\!50$, limiting statistical power for the random baseline.
(5)~We study symmetric compound sparsity ($p = s$); asymmetric configurations may yield different interaction patterns.

% ============================================================
\section{Conclusion}
\label{sec:conclusion}

We discover that compound sparsity severity is not a fixed constant but an emergent property of the strategy $\times$ phase interaction.
Importance-based KV eviction under head pruning produces monotonically increasing compound interaction severity across reasoning phases, with paired median differences spanning an order of magnitude ($0.39 \to 8.48$) and large effect sizes ($d_z$ up to $1.44$).
Cross-architecture replication ($n\!=\!250$) on LLaMA confirms the interaction is architecture-universal in direction, while cross-benchmark replication on GPQA Diamond establishes that strategy specificity ($13\times$ importance/random ratio) generalizes to graduate-level reasoning.
Discrimination collapse---the degradation of importance score selectivity after head pruning---provides a mechanistically grounded explanation consistent across all three models.
These findings motivate phase-aware compound compression as a principled alternative to uniform sparsity policies.

% ============================================================
\appendix
\section{Experimental Configuration}
\label{app:config}

\begin{table}[h]
\centering
\small
\caption{Complete experimental configuration.}
\begin{tabular}{@{}ll@{}}
\toprule
Parameter & Value \\
\midrule
Pruning rate & $p = s$ (symmetric) \\
Head selection & Lowest attention score heads \\
KV eviction (importance) & Lowest cumulative attention \\
KV eviction (random) & Uniform random \\
Phase segmentation & Quantile-based (dense model profile) \\
Bootstrap resamples & $B = 10{,}000$ (paired per-problem) \\
Bonferroni correction & $k = 4$ (sparsity levels) \\
Random seed & 42 \\
Precision & bf16 \\
\bottomrule
\end{tabular}
\end{table}

\section{Per-Phase Interaction Index on GPQA}
\label{app:gpqa_phases}

\begin{table}[h]
\centering
\small
\caption{Per-phase median interaction index on GPQA Diamond (DeepSeek, importance).}
\begin{tabular}{@{}lccc@{}}
\toprule
Sparsity & \textsc{High} & \textsc{Mid} & \textsc{Low} \\
\midrule
$s = 0.15$ & 0.012 & 0.145 & 0.528 \\
$s = 0.20$ & 0.028 & 0.418 & 1.825 \\
$s = 0.25$ & 0.065 & 1.562 & 5.842 \\
$s = 0.30$ & 0.112 & 3.245 & 12.856 \\
\bottomrule
\end{tabular}
\end{table}

\begin{table}[h]
\centering
\small
\caption{Per-phase median interaction index on GPQA Diamond (DeepSeek, random).}
\begin{tabular}{@{}lccc@{}}
\toprule
Sparsity & \textsc{High} & \textsc{Mid} & \textsc{Low} \\
\midrule
$s = 0.25$ & 0.682 & 0.845 & 1.124 \\
$s = 0.30$ & 1.856 & 2.124 & 2.812 \\
\bottomrule
\end{tabular}
\end{table}

\section{LLaMA Detailed Results}
\label{app:llama}

\begin{table}[h]
\centering
\small
\caption{LLaMA-8B ($n\!=\!250$) detailed results. Phase gradient: within-importance $d_z$ and within-random $d_z$; strategy interaction: imp-vs-rand $d_z$ with Bonferroni $p$.}
\begin{tabular}{@{}lcccc@{}}
\toprule
& $s\!=\!0.15$ & $s\!=\!0.20$ & $s\!=\!0.25$ & $s\!=\!0.30$ \\
\midrule
Imp phase $d_z$ & 1.02 & 1.09 & 1.12 & 1.05 \\
Rand phase $d_z$ & 0.36 & 0.38 & 0.42 & 0.41 \\
Imp/Rand ratio & 2.1$\times$ & 2.2$\times$ & 2.6$\times$ & 3.4$\times$ \\
\midrule
Imp-vs-rand $d_z$ & 0.24 & 0.31 & 0.38 & 0.43 \\
Bonf.\ $p$ & 0.0013 & $3.6\!\times\!10^{-5}$ & $5.6\!\times\!10^{-8}$ & $8.4\!\times\!10^{-10}$ \\
Median diff & 0.018 & 0.052 & 0.285 & 0.612 \\
\bottomrule
\end{tabular}
\end{table}

% ============================================================
\begin{thebibliography}{20}

\bibitem[{SparK}(2024)]{spark2024}
SparK: Joint Sparse Attention and KV Cache Compression for Efficient LLM Inference. \emph{Preprint}, 2024.

\bibitem[{SLiM}(2024)]{slim2024}
SLiM: Sparse and Lightweight Model Compression via Compound Sparsity. \emph{NeurIPS}, 2024.

\bibitem[{Ada-KV}(2024)]{adakv2024}
Ada-KV: Adaptive KV Cache Eviction with Per-Head Budget Allocation. \emph{EMNLP}, 2024.

\bibitem[{ThinKV}(2025)]{thinkv2025}
ThinKV: Phase-Aware KV Cache Compression for Chain-of-Thought Reasoning. \emph{arXiv:2510.01290}, 2025.

\bibitem[{R-KV}(2025)]{rkv2025}
R-KV: Reasoning-Optimized KV Cache Compression. \emph{arXiv:2505.24133}, 2025.

\bibitem[{HIES}(2025)]{hies2025}
HIES: Importance Score Collapse and Entropy-Compensated Eviction. \emph{arXiv:2510.13832}, 2025.

\bibitem[{Lethe}(2026)]{lethe2026}
Lethe: Layer-Time Adaptive KV Cache Pruning. \emph{AAAI}, 2026.

\bibitem[{Physics of KV}(2026)]{physicskv2026}
The Physics of KV Cache Compression: Phase Transitions and Representational Rigidity. \emph{arXiv:2603.01426}, 2026.

\bibitem[{GPQA}(2024)]{gpqa2024}
GPQA: A Graduate-Level Google-Proof Q\&A Benchmark. \emph{ICLR}, 2024.

\end{thebibliography}

\end{document}
